% #############################################################################
% This is Chapter 6
% !TEX root = ../main.tex
% #############################################################################
\fancychapter{Py2HWSW Layer}
\label{chap6:fatori_hardware}




% #############################################################################
\section{Fault Tolerance Support Infrastructure}
\label{chap6:starting}

The IOb-\ac{SoC}-Ibex described in \cref{chap5:iob-soc-ibex} was developed to stay as close to the native Ibex implementation as possible. In order to meet \ac{Fatori-V}'s requirements, preliminary support infrastructure was added to make the system externally configurable and easier to modify at the \ac{RTL} level.

As described in \cref{chap3:hwvssw}, \ac{Fatori-V} requires the \ac{RTL} to be configured through macro files. IOb-\ac{SoC}-Ibex was modified to use two macro files. The \texttt{fatori\_features.svh} file controls architectural features, such as the ICache, branch prediction, the Writeback Stage, and \ac{ISA} extensions. The \texttt{fatori\_ftm.svh} file controls native Ibex \acp{FTM}. In native Ibex, the SecureIbex flag enables all \acp{FTM} as one group. Through \texttt{fatori\_ftm.svh}, those mechanisms can be enabled independently.

The native Ibex clock-gating cell was disabled from the Ibex top module, and the internal clock was wired directly to \texttt{clk\_i}, from the \ac{SoC}. This allowed the implemented architecture to meet timing constraints.

Ibex originally implemented 186 pipeline registers directly with \texttt{always\_ff} blocks spread across the source files, many times inside cascaded \texttt{if} conditions or complex \acp{FSM}. This made register-level modifications harder to implement. Register Redundancy, planned since the WeeRwolf EL2 work in \cref{chap4:el2}, requires a register to be replaced by multiple replicas and a voter, which cannot be done cleanly while each register is implemented as a separate \texttt{always\_ff} block.

Each pipeline register was replaced by a \texttt{FATORI\_REG()} or \texttt{FATORI\_REG\_ENUM()} macro call, developed for \ac{Fatori-V} and defined in \texttt{fatori\_macro\_functions.svh}. These macros act as register wrapper entry points. They keep the register call sites uniform and expand to the selected register hardware implementation. By default, the macro expands to IOBundle's \texttt{iob\_reg\_re}. Other register wrappers were implemented for \ac{Fatori-V}, and are described in subsequent sections.

\texttt{FATORI\_REG()} is used for logic vectors. \texttt{FATORI\_REG\_ENUM()} is used for typed registers by packing the input data into a logic vector with \texttt{\$bits()} and casting it back on output. Each register enable logic was also reduced to one explicit \texttt{enable} per-register signal.

Each wrapped register has a unique prefix used to derive internal signal names and identifiers. The file \texttt{fatori\_registers.yaml}, introduced in \cref{chap3:fatori-v}, stores the prefix, hierarchy path, and ID of every wrapped register.






% #############################################################################
\section{Fault Tolerance Mechanisms}
\label{chap6:ftms}

Although Ibex already includes native \acp{FTM}, this work expanded that set to make \ac{Fatori-V} a broader evaluation platform. This section details the implementations of the M-of-N Redundancy and Self-Testing \acp{FTM}. 



% #############################################################################
\subsection{Hardware Redundancy}
\label{chap6:mon}

As introduced in \cref{chap2:ftms}, M-of-N Redundancy (or M-of-N Systems) replicates a block \(N\) times and accepts an output when at least \(M\) replicas agree. In \ac{Fatori-V}, this is applied to wrapped pipeline registers and to selected logic modules. When correction is enabled, the same flow also applies scrubbing by restoring replicas that disagree with the voted value.

Register M-of-N extends the register wrapper introduced in \cref{chap6:starting}. Each register's M-of-N parameters are configured through a new single macro file, \texttt{fatori\_reg\_mon.svh}. The \texttt{*\_MON\_N} macro selects the number of replicas, \texttt{*\_MON\_M} selects the quorum threshold, and \texttt{*\_MON\_CORRECTION} enables correction. When \texttt{M=0}, the threshold is set to strict majority, \(\lfloor N/2 \rfloor + 1\), by the hardware.

When redundancy is disabled, the wrapper keeps the default register implementation. When redundancy is enabled, it instantiates exactly \(N\) physical register replicas using Verilog \texttt{generate} logic. All replicas receive the same input and enable signal.

The register voter compares the replica values and counts how many replicas match each value. A value is accepted only if it reaches the configured quorum and no different value also reaches that quorum. If one or more replicas disagree with the voted value, the wrapper reports a minor error and marks the disagreeing replicas. If no unique value reaches \(M\), the wrapper reports a major error. When correction is enabled, the voted value is written back into the disagreeing replicas, restoring corrupted data.

Logic M-of-N applies the same voting principle to larger hardware modules. Registers can share one common wrapper because they have the same basic interface. Logic modules do not, so each protected target needs a wrapper adapted to its native ports. The need for individual wrappers was the main limiting factor when expanding Logic M-of-N to other modules. In this work, wrappers were developed for \texttt{ibex\_if\_stage}, \texttt{ibex\_controller}, \texttt{ibex\_decoder}, \texttt{ibex\_lsu}, and \texttt{ibex\_alu}.

Logic redundancy is configured through another macro file, \texttt{fatori\_logic\_mon.svh}. Each target uses \texttt{*\_MON\_N}, \texttt{*\_MON\_M}, and \texttt{*\_MON\_HOLD}. When \(N=1\), the native module is instantiated. When \(N>1\), the wrapper instantiates exactly \(N\) replicas with the same inputs and votes their outputs.

Logic M-of-N does not include correction, because that would require accessing the \ac{FPGA}'s configuration memory. Instead, the \texttt{HOLD} feature outputs the last valid voted value when no quorum exists.

Each logic wrapper groups the target module's outputs into a packed structure. This structure is converted into a bit vector and sent to \texttt{fatori\_mon\_voter}. The voter receives the replica outputs as an array and uses the configured \(N\) and \(M\) values to compare them. Since \(N\), \(M\), and the voted width are parameters, synthesis elaborates only the comparison and voting hardware required by the selected configuration.

The voter accepts the value produced by at least \(M\) replicas. When a valid value exists, the selected bit vector is cast back to the output structure used by the wrapper, and it is exported to the pipeline.



% -------------------------------------------------------------
\subsection{Self-Testing Hardware}
\label{chap6:self_test}

As introduced in \cref{chap2:ftms}, Self-Testing checks a hardware block by applying predefined inputs and comparing the outputs with known results. In this work, it was implemented for hardware blocks that can be tested with deterministic input patterns and whose idle periods can be identified from existing Ibex signals.

Each self-test wrapper keeps the same interface as the native block. Internally, it instantiates the native module and the developed \ac{ST} Core. The wrapper computes the \texttt{idle} signal. The \ac{ST} Core reads this signal, stores the test inputs and expected outputs, drives the native module through a multiplexer during a test, compares the produced output with the expected result, and reports mismatches.

Through the \ac{ST} Core, the Self-Testing \ac{FTM} can be configured with the \texttt{fatori\_selftest.svh} macro file. Two macros per \ac{ST} target are defined. \texttt{*\_ST\_SPACING} defines the minimum number of cycles between two tests for the same target. \texttt{*\_ST\_POLICY} defines what happens when the tested block is requested by the pipeline during a self-test. Two policies were implemented. The \texttt{abort} policy allows a test to start, but cancels it if a real operation needs the target. The \texttt{hold} policy lets the test finish by holding the pipeline when supported by the wrapper.

\ac{ST} wrappers were developed for \texttt{ibex\_alu}, \texttt{ibex\_multdiv\_*}, \texttt{ibex\_load\_store\_unit}, and the \ac{PC} adder. The \ac{ALU} tests use identity operations such as \(X+0=X\) and \(X\oplus0=X\). The multiplier and divider tests use cases such as \(0\times X=0\) and \(X/1=X\). The \ac{LSU} tests exercise alignment and byte-enable behaviour without unsafe external writes. The \ac{PC} adder tests check the expected \(+2\) and \(+4\) increments.

Test execution is controlled inside the \ac{ST} Core through three inputs connected by the wrapper. \texttt{enable} allows the \ac{ST} Core to run tests. \texttt{kick} requests an immediate test. \texttt{idle} is computed by each wrapper from native Ibex signals and tells the \ac{ST} Core that the target can be tested without interfering with the current instruction. 

Because each target has its own interface and usage pattern, each wrapper defines \texttt{idle} from the native Ibex signals that were identified as reliable indicators of whether that target is currently being used. The \ac{ALU} wrapper marks the target as idle during \ac{CSR}, branch, or jump instructions, as long as no \ac{LSU} request is active. The multiplier and divider wrapper uses the native \texttt{inflight} signal to avoid testing while a multiplication or division is active. The \ac{LSU} wrapper marks idle only when there are no active requests and no pending completions. The \ac{PC} adder wrapper marks idle when the IF stage is stalled.

A stable Self-Testing implementation was not reached within the project time. The wrappers and \ac{ST} Core structure were developed, but the mechanism was not included in the final experimental evaluation.










% #############################################################################
\section{Fault Monitoring and Control Hardware}
\label{chap6:control_report}

The \acp{FTM} described in the previous sections report errors, but those reports must still be collected, counted, and exposed to firmware. For that purpose, \ac{Fatori-V} adds monitoring and control hardware using, and extending, the native Ibex error interface described in \cref{chap4:ibex_ft}.

The extended interface keeps the minor and major error native signals and adds the one-cycle pulses \texttt{new\_min} and \texttt{new\_maj}, which mark each new detected error for counting. It also adds \texttt{scrub\_occurred}, which marks corrections performed by redundancy voters.

Error signals are merged as they propagate through the core hierarchy. Each level combines locally generated events with events received from lower levels. The merged signals are then connected to the Fault Manager, described bellow, which counts events, exposes monitor values through custom \acp{CSR}, and drives the control outputs used after major errors.




% -------------------------------------------------------------
\subsection{Fault Manager Hardware Module}
\label{chap6:fault_mgr}

The developed \texttt{fatori\_fault\_mgr} module is the main \ac{FT} control block added in this work. It receives the merged minor, major, correction, and injection event signals, and drives the control outputs used to stop instruction fetching or request a reset. Its interface includes the error inputs, \texttt{core\_sleep\_i}, \texttt{fetch\_enable\_o}, and \texttt{core\_reset\_req\_o}. Additional monitor inputs and counter outputs are included when \ac{FI} or extra monitoring hardware is enabled.

The amount of monitoring hardware is selected by the Metric Level configured in \texttt{fatori\_features.svh}. Each level keeps the hardware from the lower levels and adds the monitors listed in \cref{tab:fatori_metric_levels}.

\begin{table}[h]
\centering
\small
\setlength{\tabcolsep}{3.2pt}
\renewcommand{\arraystretch}{1.20}
\begin{tabularx}{\textwidth}{|p{0.16\textwidth}|X|}
\hline
Metric Level &
What is measured \\
\hline
\hline
0 &
Baseline Ibex counters, including the number of cycles and number of retired instructions. \\
\hline
1 &
Native Ibex \ac{HPM} counters for stalls, memory accesses, jumps, branches, compressed instructions, multiplication stalls, and division stalls. \\
\hline
2 &
Minor and major error events, plus register and logic injection counts when \ac{FI} is enabled. \\
\hline
3 &
Correction events from scrubbing, internal major errors, bus major errors, and double faults. \\
\hline
4 &
Time to first minor error, time to first major error, and last injection-to-detection latency. \\
\hline
5 &
Detection latency sum and number of latency samples. \\
\hline
\end{tabularx}
\caption{Hardware monitors for each Metric Level.}
\label{tab:fatori_metric_levels}
\end{table}

The fault manager implements fixed-width counters for the enabled Metric Level. Error counters increment from the minor and major event signals. Correction counters increment from \texttt{scrub\_occurred}. Injection counters increment from the register and logic injection pulses produced by the \ac{FI} hardware path.

To implement the timing monitors, an internal cycle counter provides the time reference. The time-to-first-event registers store the cycle of the first minor or major event after reset. Detection latency stores the cycle of the last injection and compares it with the cycle of the next detected error. At the highest Metric Level, the manager also accumulates the latency sum and the number of latency samples.

The fault manager can also define the response to major errors. When enabled, it disables instruction fetching so the core stops starting new instructions while the pipeline drains. It can then request a reset through \texttt{core\_reset\_req\_o}. Reset behaviour is configured through \texttt{FATORI\_RESET\_ON\_MAJOR} and \texttt{FATORI\_WAIT\_SLEEP\_BEFORE\_RESET}. The second macro delays the reset request until the core reports that it is idle.





% -------------------------------------------------------------
\subsection{Monitoring CSRs}
\label{chap6:metrics_hw}

The monitoring hardware reuses native Ibex performance \acp{CSR} and adds \ac{FT}-specific counters. The new monitor values are wired into \texttt{ibex\_cs\_registers} and exposed as read-only custom \acp{CSR} in the \texttt{0xBC0} address range. Firmware reads them with \texttt{csrr}.

\begin{table}[h]
\centering
\small
\setlength{\tabcolsep}{3.2pt}
\renewcommand{\arraystretch}{1.25}
\begin{tabularx}{\textwidth}{|l|l|X|}
\hline
\ac{CSR} name &
Address &
Purpose \\
\hline
\hline
\texttt{CSR\_FATORI\_MINOR\_CNT} &
\texttt{0xBC1} &
Counts minor error events. \\
\hline
\texttt{CSR\_FATORI\_MAJOR\_CNT} &
\texttt{0xBC2} &
Counts major error events. \\
\hline
\texttt{CSR\_FATORI\_REG\_INJ\_CNT} &
\texttt{0xBC3} &
Counts register injections issued by the \ac{FI} hardware path. \\
\hline
\texttt{CSR\_FATORI\_LOGIC\_INJ\_CNT} &
\texttt{0xBC4} &
Counts logic injections issued by the \ac{FI} hardware path. \\
\hline
\texttt{CSR\_FATORI\_CORRECTED\_CNT} &
\texttt{0xBC5} &
Counts corrections reported by scrubbing. \\
\hline
\texttt{CSR\_FATORI\_MAJOR\_INTERNAL\_CNT} &
\texttt{0xBC6} &
Counts major errors reported inside the core. \\
\hline
\texttt{CSR\_FATORI\_MAJOR\_BUS\_CNT} &
\texttt{0xBC7} &
Counts major errors reported by bus-related checks. \\
\hline
\texttt{CSR\_FATORI\_DOUBLE\_FAULT\_CNT} &
\texttt{0xBC8} &
Counts double-fault events. \\
\hline
\texttt{CSR\_FATORI\_CYCLES\_MIN} &
\texttt{0xBC9} &
Stores the cycle count of the first minor error. \\
\hline
\texttt{CSR\_FATORI\_CYCLES\_MAJ} &
\texttt{0xBCA} &
Stores the cycle count of the first major error. \\
\hline
\texttt{CSR\_FATORI\_DETECT\_LATENCY} &
\texttt{0xBCB} &
Stores the last injection-to-detection latency. \\
\hline
\texttt{CSR\_FATORI\_LATENCY\_SUM} &
\texttt{0xBCC} &
Accumulates detection latency values. \\
\hline
\texttt{CSR\_FATORI\_LATENCY\_CNT} &
\texttt{0xBCD} &
Counts detection latency samples. \\
\hline
\end{tabularx}
\caption{\ac{Fatori-V} monitoring \acp{CSR}.}
\label{tab:fatori_ft_csrs}
\end{table}

The amount of monitoring hardware is selected through the Metric Level configured in \texttt{fatori\_features.svh}. Level 1 exposes only native Ibex hardware-performance counters. Level 2 adds basic fault counters and, when \ac{FI} is enabled, injection counters. Level 3 adds corrected-event counting and the major-error breakdown. Level 4 adds time-to-first-event registers and last detection latency. Level 5 adds latency sum and latency count.

Each active \ac{CSR} is read by firmware through dedicated functions in \texttt{iob\_csr.h}. The enabled monitor values are sampled before and after each benchmark session. The firmware serializes these samples into a buffer, together with benchmark-specific values, and transfers the buffer over \ac{UART}. The host receives the transfer and writes the raw values into \texttt{metrics.txt}.








% #############################################################################
\section{Fault Monitoring and Control Hardware}
\label{chap6:control_report}

This section defines the error and metrics interface used by \ac{Fatori-V}. It shows how native and added \acp{FTM} signals are merged and counted in hardware, then exposed through custom \acp{CSR} for firmware and the Software Layer to read.

% -------------------------------------------------------------
\subsection{Fault Monitor Hardware}
\label{chap6:metrics_hw}

\ac{Fatori-V} extends Ibex's native Error Interface detailed in \cref{chap5:ibex_ft}, providing a 5 signal interface containing:
(i) the native Ibex \texttt{minor} and \texttt{major} error channels report a continuous system-health signal, (ii) two new pulses, \texttt{new\_min} and \texttt{new\_maj}, that pulse each new error, and (iii) a \texttt{scrub\_occurred} pulse that reports corrections performed by replication voters. The other signals of \cref{chap5:ibex_ft} still remain.

All error signals are aggregated (via the OR operator) progressively as they propagate upwards. At each hierarchy level, locally generated events (e.g., a voter disagreement) are merged with events coming from below (e.g., errors produced inside a replicated module, including those raised by wrapped registers within that module). Merging is done per clock cycle: multiple pulses of the same class in the same cycle collapse into a single pulse.

The Hardware Layer implements monitors used to evaluate \ac{FT} experiments by reusing the native Ibex performance \acp{CSR} (\cref{chap5:ibex_overview}) and extending them with \ac{FT}-specific hardware counters. The exported monitoring \acp{CSR} added in this work are listed in \cref{tab:fatori_ft_csrs}.

\begin{table}[H]
\centering
\begin{tabularx}{\textwidth}{|l|l|X|}
\hline
New \ac{CSR} name & \ac{CSR} Addr & Counts\\
\hline
\texttt{CSR\_FATORI\_MINOR\_CNT}          & \texttt{0xBC1} & minor error count\\
\hline
\texttt{CSR\_FATORI\_MAJOR\_CNT}          & \texttt{0xBC2} & major error count\\
\hline
\texttt{CSR\_FATORI\_REG\_INJ\_CNT}       & \texttt{0xBC3} & register fault injection count\\
\hline
\texttt{CSR\_FATORI\_LOGIC\_INJ\_CNT}     & \texttt{0xBC4} & logic fault injection count\\
\hline
\texttt{CSR\_FATORI\_CORRECTED\_CNT}      & \texttt{0xBC5} & corrected (scrub) count\\
\hline
\texttt{CSR\_FATORI\_MAJOR\_INTERNAL\_CNT}& \texttt{0xBC6} & internal-major error count\\
\hline
\texttt{CSR\_FATORI\_MAJOR\_BUS\_CNT}     & \texttt{0xBC7} & bus-major error count\\
\hline
\texttt{CSR\_FATORI\_DOUBLE\_FAULT\_CNT}  & \texttt{0xBC8} & double-fault count\\
\hline
\texttt{CSR\_FATORI\_CYCLES\_MIN}         & \texttt{0xBC9} & cycles to first minor\\
\hline
\texttt{CSR\_FATORI\_CYCLES\_MAJ}         & \texttt{0xBCA} & cycles to first major\\
\hline
\texttt{CSR\_FATORI\_DETECT\_LATENCY}     & \texttt{0xBCB} & last injection-to-detection latency\\
\hline
\texttt{CSR\_FATORI\_LATENCY\_SUM}        & \texttt{0xBCC} & detection-latency sum\\
\hline
\texttt{CSR\_FATORI\_LATENCY\_CNT}        & \texttt{0xBCD} & detection-latency count\\
\hline
\end{tabularx}
\caption{\ac{Fatori-V} new \acp{CSR} counters implemented in this work.}
\label{tab:fatori_ft_csrs}
\end{table}

All \ac{FT} hardware monitors are implemented in the developed Fault Manager (\cref{chap6:fault_mgr}). Aggregated error and correction pulses (\cref{chap6:error_agg}) increment counters when asserted. Injection counting uses the one-cycle injection pulses from \texttt{fi\_coms} (\cref{chap8:fi_coms}). Time-to-first-event latches capture the cycle count on the first occurrence after reset. Detection latency subtracts the last injection timestamp from the first subsequent error observation, and latencies are also accumulated into a sum and count. All monitors are exposed to firmware through custom read-only \acp{CSR} in the \texttt{0xBC0} range, read via \texttt{csrr}.

The amount of monitoring hardware included is configurable through five Metric Levels defined via \texttt{fatori\_features.svh}. Level 1 exports only native Ibex HPM counters. Level 2 adds basic fault counters and, when \ac{FI} is enabled, injection counts. Level 3 adds corrected count and the major-fault breakdown. Level 4 adds time to first event and per-injection detection latency. Level 5 adds latency statistics.

After the hardware monitors are published to \acp{CSR}, firmware uses the \texttt{csrr} instruction to read them. This is actually considered to be part of \ac{Fatori-V}'s Software Layer (see \cref{chap7:fatori_software}), but in order to keep metrics measuring and computing explained together, this section details how the hardware metrics reach the result tables.

Each \ac{CSR} in \cref{tab:fatori_ft_csrs} is read by a dedicated function in the firmware, implemented through the created \texttt{iob\_csr.h} file and gated by the active Metrics Level. The full set is serialized into a static buffer, called twice in the firmware wrapper: once before the benchmark (\texttt{[pre]} stage) and once after (\texttt{[post]} stage). The resulting buffer contains the \texttt{pre}, \texttt{post} and \texttt{benchmark} (added by each benchmark's wrapper) metrics:

\begin{lstlisting}[float=h, style=iob_python, caption={Data measured by the \ac{RTL} and written into a text file by the firmware.}]
[pre]  mcycle,0x000000000001A2B3
[pre]  fatori_minor_cnt,0
[benchmark]  coremark_size,36
[post] mcycle,0x000000000008F4C1
[post] fatori_minor_cnt,7
\end{lstlisting}

Once populated, the firmware transmits the buffer over \ac{UART} using the functions by IOBundle. The Py2HWSW console, running in the host, converts the buffer into a text file, naming it \texttt{metrics.txt}. On the host, the IOb-Console script intercepts this transfer and writes the file to disk under the \ac{FPGA} build directory. This is the only point at which \texttt{metrics.txt} is created as a physical file, the \ac{FPGA} itself never accesses the filesystem.

The \texttt{metrics.txt} file is considered a raw metrics file. The Results Stage of the Software Layer (\cref{chap7:results_stage}) will use these raw metrics as operands in the formulas defined centrally in the configuration files (see \cref{chap3:ui_config}), creating the results tables that are displayed in \cref{chap3:ui_results}.



% -------------------------------------------------------------
\subsection{Fault Manager Hardware Module}
\label{chap6:fault_mgr}

The \texttt{fatori\_fault\_mgr} hardware module designed and implemented in this work is the main \ac{FT} hardware controller in \ac{Fatori-V}. It receives the aggregated error interface described in \cref{chap5:ibex_ft,chap6:error_agg} and drives the system control outputs that affect instruction fetch and reset. Depending on the selected Metric Level (\cref{chap6:metrics_hw_impl}), it also implements the hardware monitors later published as custom \acp{CSR}. Its interface includes the aggregated error inputs, \texttt{core\_sleep\_i} to indicate when the core is idle, \texttt{fetch\_enable\_o} to freeze or enable fetch, and \texttt{core\_reset\_req\_o} to request reset. Additional monitor inputs and counter outputs are present when \ac{FI} and higher Metric Levels are enabled.

Inside the block, all incoming error signals are merged into two internal event pulses: one for minor-class events and one for major-class events. These pulses are then used as the increment and latch enables for the exported counters and timing registers, for example, they increment the minor and major counters and, on the first occurrence after reset, they latch the current cycle count to record time-to-first-event. Ibex fetch control uses a multi-bit encoded signal to reduce the chance that a single-bit upset flips the enable state. The manager drives that signal to disable fetch on major events and drain the pipeline. 

The block is configurable through build-time macros defined in \texttt{fatori\_features.svh}. Reset is controlled by a small \ac{FSM}, controlled by: \texttt{FATORI\_RESET\_ON\_MAJOR} that requests a reset after a major event, and \texttt{FATORI\_WAIT\_SLEEP\_BEFORE\_RESET} that delays that request until \texttt{core\_sleep\_i} confirms the core is idle. Other macros define the Metric Level and whether \ac{FI} is enabled. The monitors themselves are implemented as event-driven counters and latches, using the aggregated error pulses (\cref{chap6:error_agg}) and the single-cycle injection pulses from the \ac{FI} router (\cref{chap8:fi_coms}). Most monitors follow the same pattern: a pulse increments a counter, and selected events latch the current cycle count once. Measurements that depend on time use an internal cycle counter and stored timestamps, for example injection-to-detection latency is computed from the last injection timestamp and the next observed error pulse, and optional latency statistics accumulate a sum and count. All internal registers follow the same wrapping approach described in \cref{chap6:reg_wrap}. The resulting monitor values are wired into \texttt{ibex\_cs\_registers} and exposed as custom read-only \acp{CSR} in the \texttt{0xBC0} address range (\cref{chap6:metrics_hw_impl}), so firmware can read them with \texttt{csrr} and the Software Layer (\cref{chap7:fatori_software}) can build the results tables.